\subsection{Complexity Under a Dynamic Expert Registry}
\label{subsec:complexity_dynamic_registry}

\paragraph{Three computational regimes.}
The computational discussion is simplest when the \emph{parameter-estimation stage} is separated from the \emph{online routing stage}. In our notation, \(\mathcal E_t\) denotes the set of experts available at round \(t\), whereas \(\mathcal K_t\) denotes the maintained registry of experts for which the router stores expert-specific latent state. The same online routing rule can be used in three regimes:
\begin{enumerate}[leftmargin=1.5em]
    \item \textbf{Fixed-parameter online routing.} Model parameters are fixed before deployment, and the online cost is only the per-round filtering, scoring, and belief update cost.
    \item \textbf{Offline EM + online routing.} A one-time offline estimation phase is run first, after which the same online router is deployed with the learned parameters.
    \item \textbf{Offline + online EM.} Offline EM initializes the parameters and periodic online EM updates are added during deployment.
\end{enumerate}
This distinction is important for interpretation. The online decision layer is modular with respect to parameter estimation: any estimation procedure that provides the SLDS parameters can in principle be used, whereas the per-step routing complexity is determined by the filtering and selection operations once those parameters are fixed. Consequently, reviewer-facing runtime numbers should separate one-time estimation cost from steady-state online routing cost; otherwise the computational claim is ambiguous.

\paragraph{Online routing cost (no EM).}
We now describe the no-EM regime, which is the relevant online setting for the Jena and Melbourne experiments. Let \(M\) be the number of regimes, let \(d_g = \dim(g_t)\), let \(d_\alpha = \dim(u_{t,k})\), and let \(S\) denote the Monte Carlo budget used in the \(g_z\) information term. A conservative dense-linear-algebra bound for one online step is
\begin{equation}
T_t^{\mathrm{online}}
=
\widetilde O\!\left(
C_\theta(x_t)
+
M d_g^3
+
M |\mathcal K_t| d_\alpha^3
+
|\mathcal E_t|\, C_{\mathrm{score}}(M,d_g,d_\alpha,S)
\right),
\label{eq:complexity_online}
\end{equation}
where \(C_\theta(x_t)\) is the cost of evaluating the context-dependent transition model and \(C_{\mathrm{score}}\) covers predictive scoring together with the IDS-style information term. Equation~\eqref{eq:complexity_online} makes the practical dependence explicit: propagation of the latent filter scales with the maintained registry size \(|\mathcal K_t|\), whereas action scoring scales only with the currently available set size \(|\mathcal E_t|\). The corresponding maintained online state is
\begin{equation}
\mathrm{Mem}_t^{\mathrm{online}}
=
O\!\left(
M d_g^2 + M |\mathcal K_t| d_\alpha^2
\right).
\label{eq:complexity_memory}
\end{equation}
up to lower-order vectors and bookkeeping terms. Hence both memory and the expert-specific part of filtering scale with the maintained registry, not with the cumulative number of expert identities seen over the whole horizon.

\paragraph{Registry control.}
The key architectural point is that the router stores expert-specific latent state only for \(\mathcal K_t\), not for all expert identities ever encountered. Under the pruning rule in the paper, an unavailable expert is removed once it has not been selected for more than \(\Delta_{\max}\) rounds. Since only one expert is queried at each round,
\begin{equation}
|\mathcal K_t|
\le
|\mathcal E_t|+\Delta_{\max}.
\label{eq:complexity_registry_bound}
\end{equation}
Thus, if the active set size remains bounded, the maintained online state remains controlled even when the cumulative catalog of expert identities grows over time. This is the main reason to state complexity as a function of \(|\mathcal K_t|\): in dynamic-availability settings, the cumulative catalog size can grow substantially, whereas the actual online filter only needs to track the experts retained in the registry.

\paragraph{Comparison under growing expert catalogs.}
To make this distinction explicit, let
\begin{equation}
K_t^{\mathrm{cat}}
:=
\left|
\bigcup_{s=1}^{t} \mathcal E_s
\right|
\end{equation}
denote the cumulative number of expert identities encountered up to round \(t\). Under expert churn, it is possible that \(|\mathcal E_t|\) remains bounded while \(K_t^{\mathrm{cat}}\) grows steadily. In that regime, the SLDS router has the structural property that its expert-specific memory and propagation cost depend on \(|\mathcal K_t|\), not directly on \(K_t^{\mathrm{cat}}\). Combining \eqref{eq:complexity_memory} with \eqref{eq:complexity_registry_bound}, one obtains
\begin{equation}
\mathrm{Mem}_t^{\mathrm{SLDS}}
=
O\!\left(
M d_g^2 + M (|\mathcal E_t|+\Delta_{\max}) d_\alpha^2
\right),
\label{eq:complexity_memory_registry_bound}
\end{equation}
so if \(|\mathcal E_t|\) is uniformly bounded and \(\Delta_{\max}\) is fixed, then the maintained SLDS state is controlled independently of \(K_t^{\mathrm{cat}}\).

By contrast, the baselines considered in our experiments keep catalog-wide statistics. In the implemented LinUCB baseline, one ridge matrix \(A_j \in \mathbb R^{d \times d}\) and one vector \(b_j \in \mathbb R^d\) are stored for each expert identity \(j\), so
\begin{equation}
\mathrm{Mem}_t^{\mathrm{LinUCB}}
=
O\!\left(
K_t^{\mathrm{cat}} d^2
\right).
\end{equation}
For the implemented NeuralUCB baseline, the embedding parameters are shared, but the linear uncertainty statistics remain expert-specific, yielding
\begin{equation}
\mathrm{Mem}_t^{\mathrm{NeuralUCB}}
=
O\!\left(
hd + K_t^{\mathrm{cat}} h^2
\right),
\end{equation}
where \(h\) is the hidden dimension. For the shared-parameter linear baselines, the joint arm-conditioned feature dimension itself depends on the catalog size. With the current feature map,
\begin{equation}
d_{\mathrm{joint}}
=
d + K_t^{\mathrm{cat}} + K_t^{\mathrm{cat}} d,
\end{equation}
so the stored precision matrix scales as
\begin{equation}
\mathrm{Mem}_t^{\mathrm{Shared}}
=
O\!\left(
d_{\mathrm{joint}}^2
\right),
\end{equation}
whose dominant term is quadratic in the catalog-dependent feature dimension. Thus the comparison relevant to dynamic availability is not merely ``ours versus theirs'' in raw milliseconds; it is that the SLDS stores and propagates state for the maintained registry \(\mathcal K_t\), whereas the baselines studied here retain catalog-wide statistics or a catalog-dependent joint model. We emphasize that this does not rule out alternative pruned or compressed baselines in principle; rather, such pruning is not native to the baseline formulations implemented here, while it is built directly into the SLDS registry mechanism.

\paragraph{EM overhead.}
When offline EM is used, the total cost becomes
\begin{equation}
T_{\mathrm{total}}^{\mathrm{offEM}}
=
T_{\mathrm{offEM}}
+
\sum_{t=1}^{T} T_t^{\mathrm{online}},
\label{eq:complexity_offline_em}
\end{equation}
and when periodic online EM is added,
\begin{equation}
T_{\mathrm{total}}^{\mathrm{off+onEM}}
=
T_{\mathrm{offEM}}
+
\sum_{t=1}^{T} T_t^{\mathrm{online}}
+
\sum_{u\in\mathcal U} T_u^{\mathrm{onEM}},
\label{eq:complexity_online_em}
\end{equation}
where \(\mathcal U\) is the set of online adaptation times. The important practical distinction is therefore:
\begin{itemize}[leftmargin=1.5em]
    \item online runtime tables should report \(\sum_t T_t^{\mathrm{online}}\) or per-step milliseconds,
    \item any offline EM warm-up should be reported separately as a one-time training cost.
\end{itemize}
This separation is the correct interpretation of our experiments: Jena and Melbourne are reported in the fixed-parameter online-routing regime, whereas any experiment using EM should report its one-time fitting cost separately from the steady-state routing cost.

\subsection{Registry-Churn Experiment on the Paper Synthetic}
\label{subsec:complexity_registry_experiment}

\paragraph{Setup.}
We use the same synthetic residual generator as in the paper's main synthetic benchmark: the two-regime \texttt{tri\_cycle\_corr} model with partial feedback, shared latent factor, and model-matched residual generation. The only change is the availability process. We keep the active set size fixed at
\begin{equation}
|\mathcal E_t| = 4 \qquad \text{for all } t,
\end{equation}
while increasing the cumulative catalog size
\begin{equation}
R \in \{4,8,16,32,64\}.
\end{equation}
Experts are partitioned into two latent families, matching the paper's synthetic structure: two family-A slots and two family-B slots remain active at all times, and every \(30\) rounds one active slot is replaced by a fresh expert from the same family. This yields a bounded active set with growing cumulative expert identities. We keep the paper's no-EM routing configuration with \(\Delta_{\max}=250\), and we use uniform-initialized transitions so that the complexity comparison isolates online routing rather than parameter learning.

\paragraph{Measurement protocol.}
To isolate algorithmic routing cost, contexts, losses, and residuals are precomputed once and excluded from the timer. We measure only the router/baseline \texttt{select}+\texttt{update} loop. For L2D-SLDS, ``online state entries'' means the exact maintained filter state in the implementation,
\begin{equation}
\bigl(w_t,\mu_g,\Sigma_g,\{\mu_{u,k},\Sigma_{u,k}\}_{k\in\mathcal K_t}\bigr),
\end{equation}
counted as stored array entries and averaged over time. Accordingly, the quantity reported in the tables as \emph{Avg.\ state} means the time-average number of stored scalar array entries maintained online by the method during the run. For the baselines, we analogously count their allocated online statistics/weights (\(A,b\) for LinUCB and SharedLinUCB; \(A,b_{\mathrm{lin}},W_1,b_1\) for NeuralUCB).

\begin{table}[t]
\caption{Bounded active set, growing catalog on the paper synthetic. The active set is always four experts, while the cumulative catalog grows to \(R\). Here ``Avg.\ state'' denotes the time-average number of stored scalar entries in the maintained online state. L2D-SLDS keeps a small maintained registry and modest maintained online state size; the shared-parameter linear baseline must allocate a rapidly growing joint model.}
\label{tab:complexity_registry_scaling}
\centering
\scriptsize
\resizebox{\columnwidth}{!}{%
\begin{tabular}{rrrrrrr}
\toprule
\(R\) & Mean cost & \(\overline{|\mathcal K_t|}\) & \(\max |\mathcal K_t|\) & Avg.\ state (ours) & Avg.\ state (Shared) & Avg.\ state (Neural) \\
\midrule
4  & 13.25 & 4.00 & 4  & 30.0 & 90    & 1700 \\
8  & 13.88 & 4.32 & 8  & 31.3 & 306   & 2928 \\
16 & 12.88 & 4.96 & 13 & 33.8 & 1122  & 5488 \\
32 & 14.13 & 6.23 & 13 & 38.9 & 4290  & 10608 \\
64 & 12.42 & 8.80 & 13 & 49.2 & 16770 & 17440 \\
\bottomrule
\end{tabular}%
}
\end{table}

\begin{table}[t]
\caption{Largest-catalog comparison (\(R=64\), \(T=3000\)). Here ``Avg.\ state'' again denotes the time-average number of stored scalar entries in the maintained online state. Even under heavy expert churn, L2D-SLDS remains the lowest-cost method while keeping its maintained online state two to three orders of magnitude smaller than the strongest shared baseline and NeuralUCB.}
\label{tab:complexity_registry_methods}
\centering
\scriptsize
\begin{tabular}{lrrr}
\toprule
Method & Mean cost & Cum. cost & Avg.\ state \\
\midrule
L2D-SLDS     & 12.42 & 37252 & 49.2 \\
NeuralUCB    & 15.41 & 46210 & 17440 \\
LinUCB       & 18.80 & 56382 & 128 \\
SharedLinUCB & 20.67 & 61992 & 16770 \\
\bottomrule
\end{tabular}
\end{table}

\paragraph{Interpretation.}
The experiment isolates the regime reviewers asked about: bounded active set, growing cumulative expert catalog, and no offline EM. The main facts are:
\begin{enumerate}[leftmargin=1.5em]
    \item When the catalog grows from \(R=4\) to \(R=64\), the cumulative number of expert identities seen increases by a factor of \(16\), but the maintained SLDS registry stays small: \(\overline{|\mathcal K_t|}=8.80\) and \(\max_t |\mathcal K_t| = 13\) at \(R=64\).
    \item The corresponding maintained online state for L2D-SLDS grows only from \(30.0\) to \(49.2\) average array entries, whereas the shared-parameter linear baseline grows from \(90\) to \(16{,}770\) entries, and NeuralUCB reaches \(17{,}440\).
    \item L2D-SLDS keeps the best routing performance throughout the sweep. At \(R=64\), its cumulative cost is \(37{,}252\), versus \(46{,}210\) for NeuralUCB, \(56{,}382\) for LinUCB, and \(61{,}992\) for SharedLinUCB.
\end{enumerate}
This experiment therefore supports the intended complexity claim: the dynamic registry is not only a bookkeeping device; it keeps the maintained online state controlled under expert churn while preserving a large performance gap over the baselines. We report runtime separately on the real-data benchmarks, where it is more directly comparable across methods.
